% Figure from MATH 591 notes §86 (Lecture 34; not from lecture): the three kinds of points in the proof of Cartan's formula
\begin{tikzpicture}[>=Stealth, scale=1.2]
\draw[thick, fill=gray!5] (-0.4,-1.3) rectangle (7.6,1.5);
\node[font=\small] at (7.25,1.2) {$M$};
\fill[blue!12] plot[smooth cycle] coordinates {(0.6,-0.6) (2.4,-0.9) (4.4,-0.5) (4.8,0.6) (3.0,1.1) (1.0,0.9)};
\draw[blue!60!black] plot[smooth cycle] coordinates {(0.6,-0.6) (2.4,-0.9) (4.4,-0.5) (4.8,0.6) (3.0,1.1) (1.0,0.9)};
\node[blue!60!black, font=\scriptsize] at (1.35,0.45) {$\operatorname{supp}\mathbf{X}$};
\foreach \x/\y in {2.0/0.0, 2.6/0.4, 3.3/-0.2, 3.9/0.3} { \draw[->, blue!60!black] (\x,\y) -- ++(0.35,0.12); }
\fill (2.6,-0.35) circle (1.4pt) node[below, font=\scriptsize] {$p_1$};
\draw[->, thick, red!70!black] (2.6,-0.35) -- ++(0.45,0.15);
\fill (6.2,0.2) circle (1.4pt) node[below, font=\scriptsize] {$p_2$};
\fill (4.79,0.55) circle (1.4pt) node[right, font=\scriptsize] {$p_3$};
\end{tikzpicture}
